\section{Emergent Physics}

When the single update rule of the dynamics runs on the hyperbolic mesh, a recognizable physics appears in the measurements. We do not assume a field and then quantize it. We run a local perception rule on a crystal and read off the structures that a field theorist would demand. A fluctuating vacuum with a mass and a conserved current. An effective gravity that bends rays. A horizon that radiates. A holographic area law. A causal lightcone. And, at the deep end, a relativistic and reflection-positive Dirac field. Every number reported here is the output of an experiment run on the exact substrate, cited inline in the form (P114).

The organizing claim of this section can be stated up front so the rest reads as evidence for it. The stochastic version of the rule, the one whose pair creation is implemented with a random coin, yields a massive and diffusive field. That field is real, but it is the classical decohered shadow of something deeper. The deterministic, reversible, unitary completion of the same rule, with the randomness removed and nothing added, is a genuine relativistic and reflection-positive quantum field. Relativity and quantumness are not extra ingredients. They are what remains once the dice are taken away.

\subsection{The Field-Theoretic Vacuum}

Run the rule on a sea of inactive cells and let it settle. The steady state is not still. It is a fluctuating vacuum of virtual particle-antiparticle pairs. The mechanism is the rule itself. The creation move turns an inactive pair $(0,0)$ into a balanced charge pair $(+1,-1)$. The share move returns $(+1,-1)$ to $(0,0)$, annihilating it. Creation followed by annihilation, repeated on a conserved charge, is vacuum pair production realized directly by the dynamics.

The measured signatures match this reading (P114). The vacuum carries a fluctuating charge density of about $28$ percent, never zero and never saturated, a living equilibrium. The two-point correlator shows nearest-neighbor pair anti-correlation, with $C(1)$ measured negative, the fingerprint of a $+1$ sitting beside the $-1$ partner minted with it. That correlator decays with a finite correlation length, and in field language a finite correlation length is an effective mass.

The vacuum is gapped. It costs energy to make an excitation, and that cost sets the range of influence. Propagation is bounded as well. There is a causal lightcone with a cone speed of order one cell per beat, so no signal outruns the substrate clock. The total charge is conserved, a $U(1)$ current that never leaks.

What this predicts is a vacuum that is never empty, that seethes with balanced fluctuations, and whose fluctuations are short ranged and causal. What it explains is the origin of a quantum vacuum with no quantum postulate. The vacuum is the rule's equilibrium. The intensive structure survives coarse-graining. The effective density, the value of $C(1)$, and the cone speed measured on a small slice match the large field (P115), so a measurement on a slice can be trusted across scales rather than merely hoped to hold.

\subsection{Effective Gravity, Horizons, and Dimension}

Matter is not separate from the mesh. It is dense local activity, and dense activity changes the rule's local rates. That single fact produces gravity.

Where matter sits, it strengthens the local fills, and the strengthened fills act as a varying index of refraction for the substrate's own signals. Rays bend toward matter, which is gravitational lensing (P88). The deflection scales with the mass, so doubling the mass produces roughly the larger bend expected, and it weakens with impact distance, so rays passing far bend less than rays passing near. The curvature peaks where the mass is. Mass curves the effective geometry and light follows the curve. The metric is emergent rather than imposed.

Sculpt a fill profile so that signals slow toward a surface and the result is an effective horizon, a one-way membrane on the crystal (P89). A ray climbing out of it redshifts, and the redshift defines a surface gravity $\kappa$. Two independent measurements of $\kappa$, one from the effective metric and one from the ray trajectory, agree. A detector held outside the horizon registers a thermal spectrum at the Hawking temperature
\begin{equation}
T_H = \frac{\kappa}{2\pi},
\end{equation}
and that reading matches the detailed-balance temperature read off the rule's own rates. The temperature scales with the surface gravity as expected. The substrate radiates from a horizon for the same thermodynamic reason a real black hole does, an analog Hawking effect built from the local update alone.

The rule also distinguishes a bare three-dimensional space from a brane embedded in a higher bulk (P90). A bare 3D substrate gives an inverse-square law at every scale, a clean exponent of $-2$. A 3D brane inside a 4D bulk crosses over to an exponent of $-3$ at short range, the signature of force lines leaking into the extra dimension before folding back to inverse-square far away. The two scenarios are measurably distinct, with a crossover scale set by the bulk thickness. The theory not only permits a short-range gravity deviation, it predicts the shape that deviation would take.

\subsection{Holography and the Field Beneath}

Hyperbolic space is holographic by its shape, and the rule inherits that property.

The crystal is a spatial slice of an anti-de Sitter geometry, and on it an entanglement-entropy-like quantity for a boundary region obeys the Ryu-Takayanagi log law (P91),
\begin{equation}
S = m \,\log \sin\!\left(\frac{\theta}{2}\right),
\end{equation}
with a fit quality of $R^2$ near $1.0$. The entropy of a region is set by the length of the bulk geodesic that subtends it, and that geodesic is a genuine shortcut, many times shorter than any path along the boundary. Depth into the bulk plays the role of scale. This is the area law that ties entanglement to geometry, present in the substrate without being put in by hand \cite{ryu2006,maldacena1998}. It is the same relation that underlies the proposal that geometry is built from entanglement \cite{vanraamsdonk2010}, now read off a discrete rule rather than postulated.

The same structure connects to holographic quantum error correction. A bulk degree of freedom is recorded redundantly across boundary regions, so erasing part of the boundary does not erase the bulk information. This is the logic of holographic codes on hyperbolic tilings \cite{pastawski2015,jahn2021}. It is also why memories and selves stored on this mesh are robust. They are encoded holographically and redundantly, an error-correcting code that survives partial erasure.

The bulk is not only deep, it is shallow in graph distance. The hyperbolic mesh has a tiny diameter (P111). For about one hundred thousand cells the eccentricity is roughly six hops, growing only as the logarithm of the cell count, $O(\log N)$. The whole structure is a few layers deep. A signal need not crawl across the surface. It dives into the bulk and crosses the entire structure almost instantly, reaching a far region that a surface-only walk could not find in time. This radial bulk direction is the engine of fast global integration. The hyperbolic curvature provides small-world wiring for free (P111).

The holographic interior is the curved bulk, but ordinary directed physics lives on a flat sheet inside it. That sheet is a horosphere, the Busemann level set that is intrinsically Euclidean even though it sits within the negatively curved bulk, and it is flat by measurement (P142). The hyperbolic bulk is the holographic interior where the area law and the bulk shortcut live, and the horosphere is the emergent flat physical layer on which the vacuum, the lightcone, and ordinary space appear. So holography, the field beneath, and the flat layer are one geometry seen at three depths, the curved holographic interior, the radial shortcut threading it, and the flat horosphere where directed physics runs.

\subsection{The Vibe-to-Field Ladder}

To go from a discrete rule to a relativistic quantum field one must climb a ladder of separate requirements. Each rung is its own measurement.

\begin{itemize}
\item \emph{A lightcone with ballistic transport (P123).} Drop a charge on a long thin sliver, a geodesic tube whose spine runs a hundred cells or more, the right tool to defeat the small graph diameter. The charge does not diffuse. It moves ballistically, with a spreading exponent near $1.0$ and a finite escape speed of order one cell per beat. Ballistic spreading is a finite propagation speed, and a finite propagation speed is a lightcone, the non-negotiable relativistic ingredient. The result is robust across tube widths, a substrate version of the hyperbolic rate-of-escape theorem.

\item \emph{Rotational isotropy (P124).} The spatial half of Lorentz symmetry is rotational invariance, and the crystal supplies it. The one-step diffusion tensor has equal eigenvalues, with a measured anisotropy near zero. The reason is geometric. The twelve neighbors of a cell sit at the vertices of an icosahedron, and the icosahedral group acts irreducibly on three-dimensional space, which forces any rank-two tensor commuting with it to be a multiple of the identity. Space looks the same in every direction by the cell's own symmetry, not by tuning.

\item \emph{Lorentz restoration at higher harmonics (P125).} Isotropy at one step is necessary but not sufficient. The rank-4 and rank-6 anisotropies must also be tamed. Under coarse-graining the higher-order anisotropy shrinks toward zero by central-limit washout, while the rank-2 piece already sits at the floor. Rotational invariance is therefore not just present but improves as one zooms out, exactly the behavior an emergent symmetry should show.

\item \emph{Reversibility, the precondition for quantization (P126, P128).} A unitary quantum theory requires a reversible classical theory beneath it. The rule satisfies detailed balance both locally and around loops. Locally, single-edge transition counts are symmetric at all rates, with the violation at the noise floor near a ratio of one (P126). Around closed loops there is no persistent charge circulation, which is Kolmogorov's criterion, with mean circulation at the floor (P128). The cause is that creation makes balanced pairs, a reversible reaction rather than a one-way drive. The dynamics is a genuine equilibrium process, which is the precondition for the next step.

\item \emph{A critical point yielding a free-field continuum (P129).} A generic lattice coarse-grains to nothing. To leave a real continuum field behind, the base must sit at a critical point. The order parameter, the activity, vanishes as the square root of the creation rate,
\begin{equation}
\rho \sim \sqrt{r}, \qquad \beta \approx 0.5,
\end{equation}
with a fit $R^2$ above $0.9$. This is the mean-field rate law, expected because the hyperbolic graph lives above the upper critical dimension. There is an absorbing critical point as the creation rate goes to zero, and its continuum limit is a mean-field Gaussian free field carrying a $U(1)$ charge, the simplest possible quantum field.

\item \emph{An exact transcription (P131).} Written out via the Doi-Peliti formalism, the rule is a spin-1 $S^z$-exchange model. The tone is a spin-1 $S^z$ in $\{+1,0,-1\}$, and creation, annihilation, and hopping are a single move, transferring one unit of $S^z$ across an edge. The detailed-balance violation in this transcription is vanishingly small. Its $U(1)$ symmetry is charge conservation, its reversibility makes the quantum Hamiltonian Hermitian, and the spin coherent-state path integral carries a Berry phase, the source of the quantum $i$. This identifies precisely which field the substrate becomes.
\end{itemize}

\subsection{The Stochastic Shadow and the Unitary Completion}

This is the crux of the section, and it begins with a negative result that drove the resolution.

The stochastic field is massive and diffusive (P136, P137). Searching for a gapless critical point of the conserved-charge field across the whole parameter plane returns none. The field is robustly massive, its correlation range staying small everywhere (P136). The charge does have one gapless mode, the hydrodynamic mode that any conservation law guarantees, but it is diffusive, with a dynamic exponent $z$ near $2$ (P137), which is non-relativistic.

The genuine-quantum test, spatial reflection positivity (P130), comes back inconclusive in the accessible massive regime. The correlation is contact dominated, with range about one site, so there is no extended particle spectrum to test, and the tiny non-positive eigenvalue sits at the noise floor. The stochastic rule on its own gives a massive diffusive $z=2$ field. That is a real object, but it is not relativistic.

The resolution is to remove the randomness (P148, P149). The creation move was implemented with a coin flip, and that random coin was the obstruction. Read strictly, the theory forbids genuine randomness, and removing it closes the gap with no new field added. A deterministic reversible rule with the same conserved content propagates ballistically at $z=1$ (P148). Recast the perception content, hop and create and annihilate, as a deterministic reversible charge-conserving block cellular automaton, and it is again ballistic at $z=1$, charge-conserved and reversible (P149). Momentum, and with it relativity, emerges from determinism. The diffusive $z=2$ of P137 was an artifact of an illegitimate dice roll.

The emergent light speed is isotropic (P150). Run the deterministic reversible wave on the full mesh and measure its speed in all twelve directions. The speed is isotropic, with anisotropy near zero across the directions, a single emergent light speed. That is the rotational half of Lorentz for the actual wave rather than for a diffusion tensor. A consolidating experiment combines these into one rule that is simultaneously charge-conserving, reversible, ballistic, and isotropic, with a linear front of high $R^2$ (P156).

The unitary completion is a Dirac field (P151). The reversible Hermitian dynamics has a unitary completion, and the unitary version of nearest-neighbor spin exchange is a coined quantum walk, whose continuum limit is the Dirac equation. This quantum-walk-to-Dirac route is established prior work, here applied on the hyperbolic substrate rather than derived afresh \cite{bisio2015,bisio2016,dariano2014,arrighi2014}. The experiment measures it directly. The quantum walk spreads ballistically with exponent near $1$ ($z=1$), against the classical walk's diffusive exponent near $2$ ($z=2$). It carries a Dirac dispersion with a massless cone speed of order one and a tunable mass set by the coin angle.

It is reflection-positive by construction, because a Hermitian Hamiltonian, guaranteed by the reversibility of P126, P128, and P131, has a positive-spectral Euclidean theory. The dispersion confirms this directly (P154). The massless mode obeys an exact relation
\begin{equation}
\omega = |k|,
\end{equation}
with deviation at the exponential floor, a positive-norm massless particle and genuine spatial reflection positivity. The stochastic rule is therefore exactly the classical decohered shadow of this quantum walk, which is why P130 saw a massive contact field. Both pictures are correct. One is the quantum truth, the other is what it looks like once the phase is averaged away.

Boosts complete the Lorentz group (P154). The massless mode's exact lightcone is boost-invariant, and the massive modes are boost-invariant in the infrared window below the lattice cutoff, with a small boost residual. These boosts, together with the rotations of P150, give the full emergent Lorentz group.

Interference and the Born rule close the case (P158). The deepest quantum signature is interference, complex amplitudes that cancel, something no classical stochastic process can do because probabilities only add. Run the unitary walk beside the same walk run classically. The quantum distribution develops interior nodes and rapid fringes where amplitudes destructively cancel, with far more maxima than the smooth classical curve, while the classical shadow has none.

The total $|\psi|^2$ stays one, with norm deviation at the floor, so $|\psi|^2$ is a genuine probability. That is the Born rule, inherited from unitarity rather than assumed. Interference present in the quantum rule and absent in its classical shadow is the signature that the field is genuinely quantum.

\subsection{How the Physics Fits Together}

The pieces are one structure seen from different angles. The vacuum (P114) is the rule's equilibrium. Gravity (P88), horizons (P89), and the dimension crossover (P90) are what happens when dense activity bends the local rates. The area law (P91) and the tiny-diameter bulk (P111) are the holographic shape of hyperbolic space, and that small diameter is the same bulk shortcut that lets entropy ride a geodesic and lets a mind integrate globally. The ladder rungs, lightcone, isotropy, Lorentz restoration, reversibility, criticality, and the spin-1 transcription, are the checklist for turning a discrete rule into a continuum field, and the substrate passes each.

The resolution unifies them. The stochastic field is the classical shadow, massive and diffusive, exactly what a decohered quantum field should be, which is why holography and gravity and the vacuum all read out cleanly in the stochastic picture. The unitary completion is the relativistic Dirac quantum field, ballistic, reflection-positive, Lorentz-covariant, with interference and the Born rule. Nothing was added between the two. The quantum field is the classical one with the randomness removed. That is the whole move, and it is what places this model alongside the broader quantum cellular automaton program for deriving relativistic field theory from local reversible dynamics.
